\section*{PHỤ LỤC: LỆNH DEMO VÀ MÃ NGUỒN CHÍNH}
\phantomsection\addcontentsline{toc}{section}{\numberline{} PHỤ LỤC}

\subsection*{A. Lệnh chạy toàn bộ project}
\begin{lstlisting}[language=bash, caption={Chạy pipeline end-to-end}]
cd Metro_Ethernet_MPLS_Project
sudo bash code/run_all.sh
\end{lstlisting}

\subsection*{B. Lệnh kiểm tra MPLS}
\begin{lstlisting}[language=bash, caption={Kiểm tra route và gói MPLS}]
cat results/mpls_routes.txt
cat results/frr_control_plane.txt
cat results/tcpdump_mpls.txt
sudo python3 code/topology_mininet.py --mode mpls --test-mode --stp-wait 20
\end{lstlisting}

\subsection*{C. Demo thu cong khi vao Mininet CLI}
\begin{lstlisting}[language=bash, caption={Luong demo ngan gon khi giang vien yeu cau thao tac live}]
sudo python3 code/topology_mininet.py --mode mpls
# Trong Mininet CLI:
nodes
links
host1 ping -c 3 10.2.0.21
host2 traceroute -n 10.3.0.11
exit
sudo mn -c
\end{lstlisting}

\subsection*{D. Lệnh chạy từng phép đo tu backend/GUI}
\begin{lstlisting}[language=bash, caption={Chạy một phép đo theo cặp host}]
sudo python3 code/performance_test.py --action ping --source-host host1 --destination-host user1 --mode mpls
sudo python3 code/performance_test.py --action throughput --source-host host1 --destination-host user1 --mode mpls
sudo python3 code/performance_test.py --action packet-loss --source-host host1 --destination-host user1 --mode mpls
sudo -v
python3 code/gui_monitor.py
\end{lstlisting}

\subsection*{E. Trích mã cấu hình VPLS/L2VPN}
\begin{lstlisting}[language=Python, caption={Logic CE WAN và GRETAP pseudowire trong network_config.py}]
for network, gateway in VPLS_CE_ROUTES[ce_name]:
    run(ce, f"ip route replace {network} via {gateway} dev {intf}")

run(router, f"ip link add {tunnel} type gretap local {local_id} remote {neighbor_id} key 100 ttl 64")
run(router, f"ip link set {tunnel} master {bridge}")
run(router, f"ip link set {tunnel} up")
\end{lstlisting}

\subsection*{F. Trích mã đo UDP load sweep}
\begin{lstlisting}[language=Python, caption={Các mức tải UDP dùng để đo packet loss}]
UDP_LOAD_LEVELS = [5, 20, 50, 80, 120]

for bandwidth_mbps in UDP_LOAD_LEVELS:
    run_iperf3_udp(
        src,
        dst,
        dst_ip,
        port,
        bandwidth_mbps=bandwidth_mbps,
        label="UDP LOAD SWEEP",
    )
\end{lstlisting}
